\chapter{Requirements Specification}

\section{Introduction}
This chapter delineates the conceptual and practical foundations upon which \textit{Infra App} was meticulously designed. It encapsulates the systematic identification and thorough analysis of requirements that transform an abstract idea into a structured, actionable solution. The objectives established at this stage serve as a blueprint, ensuring that every subsequent development phase adheres to clearly defined goals and quality benchmarks.

The requirements specified here form the backbone of the system’s architecture, guiding both its functional capabilities and behavioral characteristics. The analysis is methodically divided into functional and non-functional dimensions, providing a comprehensive understanding of what the system is expected to accomplish and how it should operate under various conditions. Additionally, this chapter outlines user interactions and the product backlog, serving as a navigational framework for iterative development and continuous enhancement.
\clearpage

\section{Functional Requirements}
Functional requirements capture the essential operations and behaviors that \textit{Infra App} must exhibit to fulfill its intended objectives. They define the system’s operational blueprint, detailing how users interact with and derive value from the platform. The subsections below articulate the core functionalities, emphasizing the integration of intuitive design, intelligent automation, and actionable insight.

\subsection*{Dynamic Visual Workflows}
At the heart of \textit{Infra App} lies an intelligent visual environment that empowers users to model and interpret infrastructure as dynamic, interactive workflows. This graphical representation transforms abstract architectural configurations into comprehensible visual structures, facilitating an immediate grasp of interdependencies and component relationships. By presenting complex infrastructures in a tangible, visually coherent manner, engineers can proactively identify potential bottlenecks, assess risks, and make strategic, informed decisions with heightened efficiency.

\subsection*{Control Panel Integration}
The control panel serves as the central operational hub of \textit{Infra App}, seamlessly unifying disparate functionalities into a single interface. By embedding GitLab features within this environment, users can navigate repositories, execute deployments, restart services, and manage merge requests without the need to switch contexts. This consolidation of operations not only streamlines workflows but also replaces repetitive manual tasks with automated, reliable processes. The result is an acceleration of delivery cycles, enhanced team collaboration, and a notable reduction in operational friction.

\subsection*{Machine Learning Metrics Monitoring}
A critical dimension of \textit{Infra App} is the continuous monitoring of deployed machine learning models. Through interactive dashboards, engineers can observe key performance indicators in real-time, detecting deviations in accuracy or latency before they impact production. The system’s alerting mechanisms enable proactive intervention, ensuring sustained model reliability and operational consistency. By integrating this functionality, \textit{Infra App} transcends basic infrastructure management, positioning itself as an indispensable tool for maintaining the integrity and quality of AI-driven services.

Together, these functionalities coalesce to define the operational essence of \textit{Infra App}, where visualization, control, and analytical insight converge into a unified, user-centered platform.
\clearpage

\section{Non-Functional Requirements}
Non-functional requirements articulate the qualitative attributes that govern system behavior, ensuring that \textit{Infra App} is not only capable but also robust, efficient, and user-friendly. They describe the performance, reliability, scalability, and usability expectations that the platform must satisfy under varying conditions.

\subsection*{Performance}
Performance is a pivotal consideration. \textit{Infra App} is engineered to provide immediate responsiveness to user actions, delivering real-time updates without perceptible latency. Its architecture emphasizes fluidity and robustness, guaranteeing that high workloads do not compromise operational efficiency. By maintaining consistent responsiveness, the system ensures that users can interact seamlessly, supporting uninterrupted analytical and operational activities.

\subsection*{Reliability}
Reliability ensures that \textit{Infra App} functions accurately and consistently, even in the presence of external dependencies. The platform maintains stable integrations with critical services such as GitLab, Prometheus, and AWS. Operational tasks—including deployments, synchronizations, and data retrievals—are executed without interruption, with recovery mechanisms in place to safeguard continuity and data integrity. This commitment to reliability fosters user confidence and supports mission-critical operations.

\subsection*{Scalability}
Scalability underpins the system’s ability to adapt to evolving organizational demands. Its modular design allows for seamless expansion, accommodating additional services or components without disrupting existing functionalities. As infrastructure complexity grows, \textit{Infra App} dynamically allocates resources, preserving performance and ensuring operational stability across diverse workloads and data volumes.

\subsection*{Usability}
Usability prioritizes intuitive design and accessibility. The platform’s interface is logically structured, incorporating clear visual cues to facilitate effortless navigation. Tasks can be performed efficiently without requiring extensive technical expertise, enhancing productivity and overall user satisfaction. By foregrounding user experience, \textit{Infra App} achieves both technical excellence and operational elegance.

These non-functional requirements reinforce the system’s robustness and operational excellence, establishing a dependable foundation for daily use.
\clearpage

\section{Use Case Identification and Structuring}
Use cases provide a concrete representation of user interactions with \textit{Infra App}, identifying primary actors, their responsibilities, and the operational scenarios they engage in. This analysis illuminates the system’s functional landscape, supporting both design decisions and workflow optimization.

\textbf{Engineers} can:
\begin{itemize}
    \item Visualize infrastructure in real-time through dynamic, interactive workflows.
    \item Execute deployments and restart services via the centralized control panel.
    \item Monitor machine learning model performance and respond to alerts proactively.
\end{itemize}

\textbf{Project Managers} can:
\begin{itemize}
    \item Oversee system health and operational performance.
    \item Track deployment progress and ensure seamless release management.
    \item Leverage dashboards and analytics to make informed strategic decisions.
\end{itemize}

\begin{figure}[h!]
    \centering
    % Placeholder for the use case diagram
    \fbox{\rule{0pt}{6cm}\rule{0.9\textwidth}{0pt}}
    \caption{Global Use Case Diagram of Infra App}
\end{figure}
\clearpage

\section{Product Backlog}
The product backlog serves as a structured repository of features, enhancements, and priorities for \textit{Infra App}. It provides a clear overview of functionalities to be implemented, helping the team focus on delivering value and ensuring that system capabilities align with user needs. Each entry specifies a feature, associated user story, and its relative priority to guide development and track progress.

\begin{table}[h!]
\centering
\caption{Product Backlog – Infra App}
\renewcommand{\arraystretch}{1.4}
\begin{tabular}{>{\bfseries}c m{3.5cm} m{2cm} m{8cm} >{\centering\arraybackslash}m{1.5cm}}
\toprule
ID & Feature & Story ID & User Story & Priority \\ 
\midrule
1 & Dynamic Visual Workflows & 1.1 & As an engineer, I want to view infrastructure components and their relationships visually to understand system dependencies. & High \\[1ex]
2 & Control Panel Integration & 2.1 & As an engineer, I want to deploy and restart services from one interface to save time and reduce manual effort. & High \\[1ex]
3 & ML Metrics Monitoring & 3.1 & As a data engineer, I want to track ML model performance and receive alerts on anomalies to maintain accuracy. & Medium \\[1ex]
4 & GitLab Repository Search & 4.1 & As a developer, I want to filter repositories quickly and access details within the same dashboard. & Medium \\[1ex]
5 & User Management \& Access Control & 5.1 & As an admin, I want to manage user roles and permissions to ensure secure access. & Medium \\ 
\bottomrule
\end{tabular}
\end{table}

This backlog provides a strategic framework for the system’s evolution, capturing features and enhancements in a clear and organized manner. It ensures that development aligns with user expectations and system objectives, serving as both a reference and a roadmap for progress.

\section{Conclusion}
This chapter has meticulously articulated the requirements defining \textit{Infra App}, encompassing both functional and non-functional dimensions. It elucidated user interactions, captured in detailed use cases, and structured the project’s evolution through a comprehensive product backlog.

Collectively, these elements establish a firm foundation for the forthcoming stages, where architecture, design, and implementation will translate the envisioned system into a tangible, high-quality solution. \textit{Infra App} embodies a harmonious blend of technological sophistication, operational efficiency, and user-centric design.
